\chapter{Desk Companion}

\section{Overview}
This is the last and largest project, and it is the one you are most likely to keep using after
the workshop. You will turn the microcontroller into a \textbf{Desk Companion}: a clock that sets
itself from the internet, shows the weather for your town, and wakes you up in the morning.

Every part of this build is one you have already met. The OLED screen comes from Project 6, the
rotary encoder and speaker come from Project 3, and each one keeps the same pin it had there. What
is new is the software. Until now your board has been on its own, or has been a Wi-Fi network that
other devices joined. This time it \emph{joins} your network and asks other computers on the
internet for information.

You will build:
\begin{itemize}
    \item A clock that learns the correct time from the internet
    \item A weather screen showing the current temperature, conditions, and today's high and low
    \item An alarm you set with the knob, which rings through the speaker and can be snoozed
    \item A menu of screens you page through by turning the knob
    \item Settings that are written to the board's own filesystem, so your alarm is still there
          after the microcontroller has been unplugged
\end{itemize}

\begin{figure}[H]
\centering
    \maybeincludegraphics[width=.6\linewidth]{project_7/success.jpg}
    \caption{The finished Desk Companion, showing the time and the current conditions}
\end{figure}

\begin{tcolorbox}[colback=blue!5!white,colframe=blue!40!black]
    \textbf{Learning goals:}
    \begin{itemize}
        \item Join an existing Wi-Fi network in \emph{station} mode, the opposite of the access
              point you hosted in Project 5
        \item Act as an HTTP \emph{client}: request a page from someone else's server and read the
              reply
        \item Use a geocoding service to turn a city name into latitude and longitude
        \item Pull the values you need out of JSON sent by a real public web service
        \item Set the board's real-time clock from the headers of a reply you already asked for, and
              convert UTC into local time
        \item Write your own Python classes to organize a program that does several things at once
        \item Save settings to a file on the microcontroller so they survive a power cycle
        \item Write a main loop that never blocks, so the screen, the knob, and the alarm all stay
              responsive
    \end{itemize}
\end{tcolorbox}

\pagebreak

\section{Components}
\begin{components}
    \item Seeed XIAO ESP32-C3 microcontroller
    \item 0.96-inch 128$\times$64 SSD1306 I\textsuperscript{2}C OLED display
    \item KY-040 rotary encoder module
    \item Treedix 1W 8\si{\ohm} mini speaker
    \item Breadboard
    \item Jumper wires
\end{components}

\section{Directions}

\subsubsection{Remove previous components}
Before beginning, remove the MPU-6050 accelerometer and its wires from Project 6. Leave the
microcontroller seated in the breadboard. You will be reusing the OLED display, so if it is
already wired you may still find it easier to start from an empty board, because this project
seats it in a different place to make room for the encoder.

\subsection{Creating the circuit}

% Seating the microcontroller. This is identical for every project, so the
% coordinates live here rather than in any one project chapter. The diagram
% generator reads the \bbhole definitions below for every breadboard diagram
% it draws, which is why this file is the only place they appear.
\subsubsection{Attach the microcontroller to the breadboard}
Orient the microcontroller so that the pin labeled \textbf{5V} goes into hole \bbhole{mcu.5v}{H1}
(Column H, Row 1), \textbf{GPIO2} into \bbhole{mcu.gpio2}{D1}, \textbf{GPIO20} into
\bbhole{mcu.gpio20}{H7}, and \textbf{GPIO21} into \bbhole{mcu.gpio21}{D7}. Then press it down---this
takes more force than you might expect.

\begin{figure}[H]
    \centering
    \maybeincludegraphics[width=.95\linewidth]{common/microcontroller_seated_in_breadboard.png}
    \caption{So far, so good!}
\end{figure}


\subsubsection{Ground rail distribution}
Three separate peripherals need a ground connection in this project, which is more than the
microcontroller's own \textbf{GND} row can comfortably supply. As in Project 4, we solve this by
running one wire from \textbf{GND} to the breadboard's top blue ($-$) power rail. Every ground
connection after this taps the rail instead:

\begin{center}
\begin{tabular}{|l|l|l|}
    \hline
    \textbf{Connection} & \textbf{XIAO pin} & \textbf{Jumper holes} \\
    \hline
    \connectionrow{mcu-ground}{Top GND rail}{GND}{I2}{top-2}
    \hline
\end{tabular}
\end{center}

\subsubsection{KY-040 encoder placement and wiring}
Seat the encoder's five-pin header along the front edge of the board, knob facing you, so the pins
land in silkscreen order: \textbf{GND} in \bbhole{encoder.header.ground}{A10},
\textbf{+} in \bbhole{encoder.header.power}{A11},
\textbf{SW} in \bbhole{encoder.header.sw}{A12},
\textbf{DT} in \bbhole{encoder.header.dt}{A13}, and
\textbf{CLK} in \bbhole{encoder.header.clk}{A14}.
These are the same three signal pins the encoder used in Project 3, so the code reads the knob
exactly as the Safe Cracker did:

\begin{center}
\begin{tabular}{|l|l|l|}
    \hline
    \textbf{KY-040 pin} & \textbf{Source / Pin} & \textbf{Jumper holes} \\
    \hline
    Encoder ground & Top GND rail & \bbhole{jumper.encoder-ground.start}{top-10} $\rightarrow$ \bbhole{jumper.encoder-ground.end}{E10} \\
    \connectionrow{encoder-power}{+ (VCC)}{3V3}{I3}{E11}
    \connectionrow{encoder-sw}{SW}{GPIO 8}{I6}{E12}
    \connectionrow{encoder-dt}{DT}{GPIO 9}{I5}{E13}
    \connectionrow{encoder-clk}{CLK}{GPIO 10}{I4}{E14}
    \hline
\end{tabular}
\end{center}

\subsubsection{OLED display placement and wiring}
Seat the OLED's four-pin header in column \textbf{A} further along the board, clear of the
encoder's body: \textbf{GND} into \bbhole{oled.header.ground}{A19},
\textbf{VCC} into \bbhole{oled.header.power}{A20},
\textbf{SCL} into \bbhole{oled.header.scl}{A21}, and
\textbf{SDA} into \bbhole{oled.header.sda}{A22}.

\begin{tcolorbox}[colback=yellow!10!white,colframe=yellow!50!black]
    \textbf{Check your OLED pins again:} as in Project 6, most modules read \textbf{GND, VCC, SCL,
    SDA}, but some swap \textbf{VCC} and \textbf{GND}. Read the silkscreen on your own board before
    powering it on, and power it from \textbf{3.3V} rather than 5V.
\end{tcolorbox}

\begin{center}
\begin{tabular}{|l|l|l|}
    \hline
    \textbf{OLED pin} & \textbf{Source / Pin} & \textbf{Jumper holes} \\
    \hline
    OLED ground & Top GND rail & \bbhole{jumper.oled-ground.start}{top-19} $\rightarrow$ \bbhole{jumper.oled-ground.end}{E19} \\
    \connectionrow{oled-power}{VCC}{3V3}{J3}{E20}
    \connectionrow{oled-scl}{SCL}{GPIO 7}{B6}{E21}
    \connectionrow{oled-sda}{SDA}{GPIO 6}{B5}{E22}
    \hline
\end{tabular}
\end{center}

The screen keeps the I\textsuperscript{2}C pins it used in Project 6: \textbf{GPIO 7} for the clock
line and \textbf{GPIO 6} for the data line.

\subsubsection{Speaker placement and wiring}
Insert the speaker's two Dupont leads into the front half of the board, past the display: the
positive red lead into \bbhole{speaker.lead.signal}{C26} and the negative black lead into
\bbhole{speaker.lead.ground}{C28}. The speaker returns to \textbf{GPIO 20}, the pin it used in
Project 3:

\begin{center}
\begin{tabular}{|l|l|l|}
    \hline
    \textbf{Speaker lead} & \textbf{Source / Pin} & \textbf{Jumper holes} \\
    \hline
    \connectionrow{speaker-signal}{Signal (+)}{GPIO 20}{I7}{E26}
    Speaker ground ($-$) & Top GND rail & \bbhole{jumper.speaker-ground.start}{top-28} $\rightarrow$ \bbhole{jumper.speaker-ground.end}{E28} \\
    \hline
\end{tabular}
\end{center}

\begin{figure}[H]
    \centering
    \maybeincludegraphics[width=.95\linewidth]{project_7/wiring.png}
    \caption{Complete wiring diagram for Project 7: Desk Companion.}
\end{figure}

\subsection{Programming the microcontroller}

Upload the following files to your microcontroller using the IDE:
\begin{itemize}
    \item \texttt{project\_7\_desk\_companion.py}
    \item \texttt{lib/ssd1306.py}, \texttt{lib/rotary.py}, \texttt{lib/rotary\_irq.py}, and
          \texttt{lib/debounced\_button.py} (all in the \texttt{lib/} directory)
\end{itemize}

\begin{tcolorbox}[colback=green!5!white,colframe=green!40!black]
    \textbf{Edit these three lines before you run it.} Open
    \texttt{project\_7\_desk\_companion.py} and find the block near the top marked
    \emph{Things you should change}:
    \begin{itemize}
        \item \texttt{WIFI\_SSID} and \texttt{WIFI\_PASSWORD} --- the network to join. Your
              instructor will give you the name and password to use in the classroom.
        \item \texttt{CITY} --- the place you want the forecast for. Include the state, province,
              or country when more than one place has the same name: for example,
              \texttt{"Pittsburgh, PA"}, \texttt{"Paris, France"}, or
              \texttt{"Springfield, Illinois"}.
    \end{itemize}
\end{tcolorbox}

Then open \texttt{project\_7\_desk\_companion.py} in the IDE and run it. Refer back to
Chapter \ref{ide} for IDE instructions.

The serial console reports each step as the board starts up:
\begin{verbatim}
--- Starting Project 7: Desk Companion ---
I2C bus scanned. Found addresses: ['0x3c']
No saved settings yet; using defaults.
Connected to YWIT-Workshop as 192.168.1.42
Matched location: Pittsburgh, Pennsylvania, United States (40.4406, -79.9959)
Clock set from the weather server; UTC is (2026, 9, 20, 15, 30, 4, 5, 263)
Weather: {'temperature': 69.9, 'description': 'CLOUDY', 'high': 75.0, ...}
\end{verbatim}

\subsection{Using the Desk Companion}
Turning the knob moves between screens. Pressing it selects something on the screen you are
looking at.

\begin{center}
\begin{tabular}{|l|l|}
    \hline
    \textbf{Screen} & \textbf{What it shows} \\
    \hline
    Clock    & The time in large digits, the date, and a one-line weather summary \\
    Weather  & Temperature, conditions, today's high and low, and how old the reading is \\
    Alarm    & Press to edit the hour, press again for the minute, again to arm it \\
    Set Time & Set the clock by hand, for when there is no Wi-Fi to get the time from \\
    Status   & Whether the network is up, the board's IP address, and how long it has been on \\
    \hline
\end{tabular}
\end{center}

When the alarm goes off, the screen flashes and the speaker beeps. \textbf{Press} the knob to
switch it off, or \textbf{turn} it to snooze for nine more minutes.

\subsection{How the program works}

\subsubsection{Joining a network instead of hosting one}
In Project 5 the board created its own Wi-Fi network using \texttt{network.AP\_IF}, the
\emph{access point} interface. Here it uses \texttt{network.STA\_IF}, the \emph{station} interface,
which is what your laptop and phone use to join a network somebody else is running:

\begin{lstlisting}[language=Python, caption=Joining an existing Wi-Fi network]
self.wlan = network.WLAN(network.STA_IF)
self.wlan.active(True)
if not self.wlan.isconnected():
    self.wlan.connect(WIFI_SSID, WIFI_PASSWORD)
    for _ in range(100):  # up to ten seconds
        if self.wlan.isconnected():
            break
        time.sleep_ms(100)
\end{lstlisting}

Joining is not instant, so the code checks \texttt{isconnected()} in a loop instead of assuming it
worked. If ten seconds pass without success, the program carries on as an offline clock rather than
stopping with an error.

\subsubsection{Asking another computer for a web page}
Project 5 made your board a web \emph{server}: your phone opened a socket to it and asked for a
page. This project reverses the roles. The board is the \emph{client} now, and it opens a socket to
a weather service and sends the same kind of request a browser would:

\begin{lstlisting}[language=Python, caption=An HTTP GET request built by hand]
address = socket.getaddrinfo(host, 80)[0][-1]
sock = socket.socket()
sock.settimeout(timeout)
sock.connect(address)
request = f"GET {path} HTTP/1.0\r\nHost: {host}\r\nConnection: close\r\n\r\n"
sock.send(request.encode())
\end{lstlisting}

\texttt{getaddrinfo} is the step that turns a name such as \texttt{api.open-meteo.com} into an IP
address, which is what the socket actually needs. The reply arrives as one long stretch of bytes
with the headers first, then a blank line, then the content:

\begin{lstlisting}[language=Python, caption=Splitting the reply into headers and content]
headers, _, body = raw.partition(b"\r\n\r\n")
# The first line looks like "HTTP/1.0 200 OK"; the middle word is the code.
status = int(headers.split(b" ")[1])
return status, headers, body
\end{lstlisting}

A status of \texttt{200} means the request worked. Anything else is the server telling you it could
not help, and the program prints what it said and keeps the previous reading on screen.

\subsubsection{Turning a city name into coordinates}
The forecast service still needs latitude and longitude, but students should not have to look those
numbers up themselves. Open-Meteo has a second service called a \textbf{geocoder}. The program first
sends the value of \texttt{CITY} to that service and asks for its best match:

\begin{lstlisting}[language=Python, caption=Building the city-search request]
def geocoding_path():
    return f"/v1/search?name={url_encode(CITY)}&count=1&language=en&format=json"
\end{lstlisting}

A URL cannot contain a literal space, comma, or accented letter. \texttt{url\_encode} converts the
city's UTF-8 bytes into \emph{percent encoding}: a space becomes \texttt{\%20}, a comma becomes
\texttt{\%2C}, and so on. That means a value such as \texttt{"Pittsburgh, PA"} can safely be placed
inside an HTTP request.

The geocoding reply contains a list of possible places. Asking for \texttt{count=1} keeps only the
best match. The program takes that result's coordinates and descriptive fields:

\begin{lstlisting}[language=Python, caption=Remembering the matched place]
match = results[0]
self.location = {
    "query": CITY,
    "name": match["name"],
    "admin1": match.get("admin1", ""),
    "country": match.get("country", ""),
    "latitude": match["latitude"],
    "longitude": match["longitude"],
}
\end{lstlisting}

It prints the full match in the serial console and shows the matched city at the top of the Weather
screen, so you can check that a shared name such as Springfield did not select the wrong place. The
match is also saved in \texttt{desk\_companion.json}. On later starts the board reuses those
coordinates instead of asking the geocoder again. If you change \texttt{CITY}, the saved query no
longer matches and the board automatically looks up the new place.

\subsubsection{Reading real data out of JSON}
The weather service replies with JSON, the same format Project 5's web page used to talk to the
board. \texttt{ujson.loads} turns those bytes into ordinary Python dictionaries and lists, and from
there it is plain indexing:

\begin{lstlisting}[language=Python, caption=Pulling the useful numbers out of the reply]
payload = ujson.loads(body)
current = payload["current"]
daily = payload["daily"]
self.weather = {
    "temperature": current["temperature_2m"],
    "description": describe_weather_code(current["weather_code"]),
    "high": daily["temperature_2m_max"][0],
    "low": daily["temperature_2m_min"][0],
    "fetched_ms": time.ticks_ms(),
}
\end{lstlisting}

Conditions arrive as a number rather than words: the service uses standard WMO weather codes, where
0 is a clear sky, 3 is cloudy, and 95 is a thunderstorm. The \texttt{WEATHER\_CODES} dictionary near
the top of the file turns those numbers into the short labels that fit on the screen.

\subsubsection{Knowing what time it is}
A microcontroller has no idea what time it is when it powers on, and this one never asks a separate
time service. It does not need to. \emph{Every} HTTP reply carries a \texttt{Date} header giving the
server's own clock, to the second, always in UTC and always in the same shape:

\begin{lstlisting}[caption=The start of a real reply from the weather service]
HTTP/1.1 200 OK
Date: Sun, 20 Sep 2026 15:17:37 GMT
Content-Type: application/json
\end{lstlisting}

The board already asks that server for the weather every fifteen minutes, so the same reply sets the
clock. A helper called \texttt{parse\_http\_date} looks through the header lines for that one, and
returns the six numbers a clock needs:

\begin{lstlisting}[language=Python, caption=Taking the time from the weather reply]
server_utc = parse_http_date(headers)
if server_utc is not None:
    self.set_clock_utc(server_utc)
\end{lstlisting}

Getting two useful things out of one request is worth noticing. A separate time service would be a
second thing to wait for, a second thing to get blocked, and a second thing to explain when it
fails.

The clock keeps \textbf{UTC}, the worldwide reference time, which is four or five hours ahead of
the eastern United States depending on daylight saving. Rather than make you edit that number twice
a year, the program takes it from the weather reply, which reports the offset for the coordinates
found by the city search:

\begin{lstlisting}[language=Python, caption=Turning UTC into local time]
self.utc_offset_seconds = payload["utc_offset_seconds"]

def local_time(self):
    return time.localtime(time.time() + self.utc_offset_seconds)
\end{lstlisting}

\subsubsection{Classes that each own one screen}
This program does more at once than any earlier one. Written as a single loop of
\texttt{if}-statements it would be almost impossible to follow, so instead each screen is its own
\textbf{class}: a bundle of data and the functions that work on it. Every screen knows how to draw
itself and what the knob and button should do while it is showing.

\begin{lstlisting}[language=Python, caption=The shared shape of every screen]
class Screen:
    title = "SCREEN"

    def __init__(self, app):
        self.app = app

    def draw(self, oled):
        raise NotImplementedError

    def turn(self, steps):
        """Knob moved by the given number of clicks while editing."""

    def press(self):
        """Knob button pressed."""
\end{lstlisting}

There are five screens, and each one is declared the same way. Writing
\texttt{class ClockScreen(Screen):} means ``this is a \texttt{Screen}, plus the following
differences,'' so a screen only has to describe what makes it different. Adding a sixth screen means
writing one more small class and adding it to the list; nothing else in the program changes.

Two of the screens go a step further. Setting the alarm and setting the clock by hand need exactly
the same knob behaviour: press once to start on the hour, press again for the minute, and turn to
change whichever one is underlined. Rather than write that twice, a \texttt{TimeEditScreen} class
holds the shared part, and the two screens say what is different about them --- where the hour and
minute are kept, and what to do once you have finished editing. The alarm screen also adds a third
field for arming it, which is one extra entry in a list and one extra branch in \texttt{turn}.

\subsubsection{Remembering settings after a power cut}
Everything a program keeps in a variable is gone the moment the board loses power. The
microcontroller has a small filesystem of its own, though, and it is the same filesystem your
project files live on, so the alarm and matched city can simply be written to a file:

\begin{lstlisting}[language=Python, caption=Saving and loading settings]
def save_settings(settings):
    with open(SETTINGS_FILE, "w") as handle:
        handle.write(ujson.dumps(settings))


def load_settings():
    settings = dict(DEFAULT_SETTINGS)
    try:
        with open(SETTINGS_FILE, "r") as handle:
            saved = ujson.loads(handle.read())
    except (OSError, ValueError):
        print("No saved settings yet; using defaults.")
        return settings
\end{lstlisting}

The first run has no file to read, and a file can always be damaged or hand-edited into nonsense, so
the load starts from a copy of the defaults and only replaces the keys it recognizes. A program that
reads files has to assume the file might be missing or wrong. The cached location also records the
original \texttt{CITY} query; if that line changes, the program knows the old coordinates should be
discarded and looked up again.

\subsubsection{A loop that never waits}
Project 3 played tones with \texttt{time.sleep\_ms()}, which is fine when the only job is making a
noise. An alarm clock cannot do that: while it beeps it still has to redraw the screen and watch for
you pressing the knob. So the beeping is broken into ``switch the sound on, and note when to switch
it off again'':

\begin{lstlisting}[language=Python, caption=Beeping without blocking]
def tick(self, now_ms):
    """Called every time round the main loop to keep the beeps going."""
    if not self.ringing:
        return
    if time.ticks_diff(now_ms, self.next_change_ms) < 0:
        return
    if self.sounding:
        self._off()
        self.next_change_ms = time.ticks_add(now_ms, BEEP_OFF_MS)
    else:
        self._on(ALARM_TONE_HZ)
        self.next_change_ms = time.ticks_add(now_ms, BEEP_ON_MS)
\end{lstlisting}

The main loop uses the same idea for everything slow. Rather than counting how long ago something
happened, each job stores the time it is next \emph{due}:

\begin{lstlisting}[language=Python, caption=The main loop]
while True:
    now_ms = time.ticks_ms()

    self.handle_input()
    self.check_alarm(now_ms)
    self.beeper.tick(now_ms)
    self.service_network(now_ms)

    if time.ticks_diff(now_ms, last_draw_ms) > DRAW_INTERVAL_MS:
        self.draw()
        last_draw_ms = now_ms

    time.sleep_ms(10)
\end{lstlisting}

The screen is redrawn four times a second and the weather is refetched every fifteen minutes, which
also nudges the clock back into line, all without any part of the program waiting for another.

\subsubsection{Drawing big digits}
The OLED library can only print small 8-pixel text, which is far too fine to read across a room. The
large clock digits are drawn the way a clock radio does it, out of seven rectangles per digit, with a
dictionary saying which segments each numeral lights:

\begin{lstlisting}[language=Python, caption=Seven-segment digits]
SEGMENTS = {
    "0": "ABCDEF",
    "1": "BC",
    "2": "ABDEG",
    ...
}

for name in SEGMENTS.get(character, ""):
    oled.fill_rect(*segments[name], 1)
\end{lstlisting}

\subsection{Hardware test checklist}
\begin{enumerate}
    \item On power-up the screen shows \textbf{JOINING} and then the clock face. If it shows
          \textbf{NO WI-FI}, check the SSID and password you typed.
    \item Turn the knob one click: the screen changes and you hear a small tick from the speaker.
    \item Page to the Weather screen. A temperature and condition should appear within a few
          seconds of starting up.
    \item Page to the Alarm screen, press, and set the alarm a minute or two ahead. Press through to
          \textbf{ARMED}, then watch for the bell appearing in the corner of the clock screen.
    \item When the alarm rings, press to silence it. Unplug the board, plug it back in, and check
          that your alarm time is still set.
\end{enumerate}

\begin{tcolorbox}[colback=yellow!10!white,colframe=yellow!50!black]
    \textbf{If something does not work:}
    \begin{itemize}
        \item \textbf{Screen stays dark:} Check the OLED power and ground (\bbref{jumper.oled-power.start} to
              \bbref{jumper.oled-power.end}, and \bbref{jumper.oled-ground.start} to \bbref{jumper.oled-ground.end}),
              and check that the ground rail wire \bbref{jumper.mcu-ground.start} to \bbref{jumper.mcu-ground.end} is in place.
              If the serial output shows an empty I\textsuperscript{2}C scan, check SCL on GPIO 7
              (\bbref{jumper.oled-scl.start} to \bbref{jumper.oled-scl.end}) and SDA on GPIO 6
              (\bbref{jumper.oled-sda.start} to \bbref{jumper.oled-sda.end}).
        \item \textbf{It says NO WI-FI:} Check \texttt{WIFI\_SSID} and \texttt{WIFI\_PASSWORD} for typos. The
              ESP32-C3 only joins 2.4\,GHz networks, so a 5\,GHz-only network will never appear. Networks
              that make you accept terms in a browser first cannot be used at all.
        \item \textbf{Clock is right but weather says NO DATA:} The board reached the network but not the
              weather service. Watch the serial output: it prints the reason. A free public service is
              sometimes simply overloaded, and the board retries every minute.
        \item \textbf{It says No location found:} Make \texttt{CITY} more specific and check its spelling.
              Use a full qualifier such as \texttt{"Springfield, Illinois"} rather than an abbreviation
              if the shorter search does not find the intended place.
        \item \textbf{The wrong city appears on the Weather screen:} Add the state, province, or country
              after a comma in \texttt{CITY}. Restarting automatically replaces the cached coordinates
              because the query changed.
        \item \textbf{Time is off by a whole number of hours:} The offset comes from the weather reply, so
              until the first successful fetch the clock uses \texttt{FALLBACK\_UTC\_OFFSET\_HOURS}. Check
              that the city shown on the Weather screen is the place you intended.
        \item \textbf{Knob turns the wrong way:} Change \texttt{reverse=True} to \texttt{reverse=False} in the
              \texttt{RotaryIRQ} setup, as in Project 3.
        \item \textbf{No sound:} Check the speaker signal on GPIO 20 (\bbref{jumper.speaker-signal.start} to
              \bbref{jumper.speaker-signal.end}) and its ground (\bbref{jumper.speaker-ground.start} to
              \bbref{jumper.speaker-ground.end}).
    \end{itemize}
\end{tcolorbox}

\section{Review}
In this project you built a complete, useful device that combines everything the earlier projects
introduced with a set of new software ideas. You learned:
\begin{itemize}
    \item How joining a network in station mode differs from hosting one as an access point
    \item How an HTTP request is built and how its reply is split into a status, headers, and content
    \item How geocoding turns a human-readable city name into forecast coordinates
    \item How to read real data out of a public web service's JSON
    \item How a board with no clock of its own learns the time from a reply it already needed, and how
          UTC becomes local time
    \item How writing your own classes keeps a large program readable
    \item How to store settings in a file so they outlive a power cut
    \item Why a program that has several jobs must never sit inside \texttt{sleep}
\end{itemize}

\section{Possible Extensions}
\begin{itemize}
    \item \textbf{Auto-dimming:} Reconnect the LDR and 10\,k\si{\ohm} resistor from Project 1 to GPIO 2 and use
          \texttt{oled.contrast()} to dim the display in a dark room.
    \item \textbf{Sunrise alarm:} Reconnect the WS2812B strip from Project 4 and fade it up from red to white
          over the minute before the alarm rings.
    \item \textbf{A nicer alarm:} Replace the single beep with a melody, reusing the note tables from Project 4.
    \item \textbf{A week's forecast:} Ask the service for \texttt{forecast\_days=5} and add a screen that pages
          through the coming days.
    \item \textbf{A second city:} Keep a list of city names in the settings file and let the knob choose which
          place the weather screen shows, geocoding and caching each one as it is first used.
    \item \textbf{Timer mode:} Add a screen that counts down a number of minutes you dial in, so the same
          hardware doubles as a study or cooking timer.
    \item \textbf{Remember more:} Save which screen you were on, and reopen it after a reboot.
\end{itemize}
